\section{Implemented Import Slices}

The first implemented V.9 adapters are intentionally small. The statement CSV
command is:

\begin{verbatim}
rcount import-statement-csv <csv> <output-dir>
\end{verbatim}

The CSV uses one row per contest, reporting unit, selection, and status. It has
four column families:

\begin{itemize}
  \item Contest columns: id, title, and vote-for cardinality.
  \item Selection columns: id, label, and candidate or write-in-bucket kind.
  \item Unit columns: reporting-unit id, reporting-unit kind, and parent
    jurisdiction.
  \item Count columns: status, selection votes, undervotes, overvotes, blanks,
    and counted ballots.
\end{itemize}

The adapter stores the original statement file as source evidence and gives it a
stable statement-CSV source label. It also produces normal contest,
reporting-unit, and summary NDJSON files. A package created this way is verified
by \texttt{rcount verify}; the source-hash check is reported as the same source
hash equation used by synthetic fixtures.

The second command imports a small NIST ERR/CDF-style JSON fixture:

\begin{verbatim}
rcount import-nist-cdf-json <json> <output-dir>
\end{verbatim}

That adapter reads CDF-shaped group units, contests, contest selections, vote
counts, other-count residuals, and result status fields. It maps recognized
group-unit types to RCOUNT reporting-unit kinds, preserves the source file under
a stable CDF source label, and computes residual-aware counted ballots from
selection votes, undervotes, overvotes, and blank contests.
